% !TEX program = xelatex
\documentclass[11pt,a4paper]{article}
\usepackage{algo-preamble}

\title{\textbf{L15 \textperiodcentered{} Δυναμικός Προγραμματισμός II}\\[2pt]
\large Σακίδιο \& Μέγιστη Κοινή Υπακολουθία}
\author{Algorithms Class Hub}
\date{\small Αλγόριθμοι και Πολυπλοκότητα (K17) \textperiodcentered{} ΕΚΠΑ}

\begin{document}
\maketitle

\begin{center}\itshape\small
Το πρόβλημα του σακιδίου (0/1 knapsack) με δισδιάστατο DP και ψευδοπολυωνυμική
πολυπλοκότητα, και η μέγιστη κοινή υπακολουθία (LCS) δύο συμβολοσειρών.
\end{center}

\section{Τι θα δούμε}

Δύο κλασικά DP προβλήματα --- και ένα \textbf{νέο} δίδαγμα. Στο προηγούμενο
μάθημα τα υποπροβλήματα παραμετροποιήθηκαν με \textbf{μία} μεταβλητή ($j = $
πρώτα $j$ αιτήματα). Εδώ θα δούμε ότι αυτό \textbf{δεν αρκεί πάντα}: το πρόβλημα
του \textbf{σακιδίου} χρειάζεται \textbf{δύο} μεταβλητές. Μετά, η \textbf{μέγιστη
κοινή υπακολουθία} δείχνει DP πάνω σε δύο συμβολοσειρές.

\section{Το Πρόβλημα του Σακιδίου (Knapsack)}

\subsection{Διατύπωση}

\textbf{Είσοδος:} $n$ αντικείμενα και ένα «σακίδιο». Το αντικείμενο $i$ έχει
\textbf{βάρος} $w_i > 0$ και \textbf{αξία} $v_i > 0$. Το σακίδιο αντέχει το πολύ
$W$ κιλά.

\textbf{Στόχος:} διάλεξε υποσύνολο αντικειμένων που \textbf{χωράει}
($\sum w_i \le W$) και \textbf{μεγιστοποιεί} τη συνολική αξία.

\begin{warningbox}
\textbf{Ο άπληστος αλγόριθμος αποτυγχάνει.} Φυσικό κριτήριο: διάλεξε κάθε φορά το
αντικείμενο με τον μεγαλύτερο λόγο \textbf{αξίας προς βάρος} $v_i / w_i$. Αλλά
αυτός ο λόγος αγνοεί τη χωρητικότητα: ένα αντικείμενο με τέλειο λόγο μπορεί να
«σπαταλήσει» χώρο που θα έδινε καλύτερο \textbf{συνδυασμό}. (Παράδειγμα: τα
αντικείμενα 3, 4 δίνουν αξία 40 με βάρος 11, ενώ τα 5, 2, 1 δίνουν αξία 35.)
\end{warningbox}

\subsection{Πρώτη απόπειρα --- και γιατί αποτυγχάνει}

Δοκιμάζουμε το μοτίβο του προηγούμενου μαθήματος: $\text{OPT}(i) = $ μέγιστη αξία
υποσυνόλων από τα αντικείμενα $1, \dots, i$.
\begin{itemize}
  \item \textbf{Το $i$ ΕΞΩ:} $\text{OPT}(i) = \text{OPT}(i-1)$.
  \item \textbf{Το $i$ ΜΕΣΑ:} η επιλογή του $i$ δεν αποκλείει άμεσα κανένα άλλο
  αντικείμενο\dots αλλά \textbf{καταναλώνει χώρο}. Χωρίς να ξέρουμε πόσος χώρος
  έχει ήδη ξοδευτεί από τα προηγούμενα αντικείμενα, \textbf{δεν μπορούμε καν να
  ξέρουμε αν χωράει το $i$}.
\end{itemize}

\begin{keypoint}
Η παράμετρος $i$ από μόνη της \textbf{δεν περιγράφει πλήρως ένα υποπρόβλημα}.
Χρειαζόμαστε να θυμόμαστε και \textbf{πόση χωρητικότητα απομένει}. Άρα: μία
δεύτερη μεταβλητή.
\end{keypoint}

\subsection{Η σωστή αναδρομή --- δύο μεταβλητές}

\begin{keypoint}
$\text{OPT}(i, w) = $ μέγιστη αξία υποσυνόλου των αντικειμένων $1, \dots, i$ με
\textbf{συνολικό βάρος το πολύ $w$}.
\end{keypoint}

Σκεφτόμαστε ξανά το αντικείμενο $i$:
\begin{itemize}
  \item \textbf{Το $i$ ΕΞΩ:} καλύτερη λύση από τα $1, \dots, i-1$ με
  χωρητικότητα $w$ $\to \text{OPT}(i-1, w)$.
  \item \textbf{Το $i$ ΜΕΣΑ} (μόνο αν $w_i \le w$): «ξοδεύουμε» $w_i$ χώρο και
  $v_i$ αξία, και λύνουμε το υπόλοιπο με χωρητικότητα $w - w_i$ $\to v_i +
  \text{OPT}(i-1, w - w_i)$.
\end{itemize}

\begin{keypoint}
\[
\text{OPT}(i, w) =
\begin{cases}
0 & i = 0 \\[4pt]
\text{OPT}(i-1, w) & w_i > w \\[4pt]
\max\bigl\{\, \text{OPT}(i-1, w),\ \ v_i + \text{OPT}(i-1, w - w_i) \,\bigr\} &
\text{αλλιώς}
\end{cases}
\]
\end{keypoint}

Κάθε κελί κοιτάει μόνο την προηγούμενη γραμμή: «πάνω» (το $i$ έξω) και
«πάνω-αριστερά κατά $w_i$» (το $i$ μέσα).

\subsection{Ο αλγόριθμος}

Συμπληρώνουμε έναν πίνακα $n \times W$ --- γραμμή προς γραμμή:

\begin{codebox}
\begin{verbatim}
Σακίδιο(n, w[], v[], W):
  Για w = 0 έως W:  M[0, w] <- 0
  Για i = 1 έως n:
    Για w = 0 έως W:
      Εάν w_i > w:
        M[i, w] <- M[i-1, w]
      Αλλιώς:
        M[i, w] <- max{ v_i + M[i-1, w - w_i],  M[i-1, w] }
  επίστρεψε M[n, W]
\end{verbatim}
\end{codebox}

\textbf{Παράδειγμα.} Αντικείμενα $(w, v)$: $1 = (2,3),\ 2 = (3,4),\ 3 = (4,5),\
4 = (5,6)$, χωρητικότητα $W = 8$.

\begin{center}
\begin{tabular}{@{}l *{9}{c}@{}}
\toprule
$i \backslash w$ & 0 & 1 & 2 & 3 & 4 & 5 & 6 & 7 & 8 \\
\midrule
\textbf{0}            & 0 & 0 & 0 & 0 & 0 & 0 & 0 & 0 & 0 \\
\textbf{1} $(2,3)$    & 0 & 0 & 3 & 3 & 3 & 3 & 3 & 3 & 3 \\
\textbf{2} $(3,4)$    & 0 & 0 & 3 & 4 & 4 & 7 & 7 & 7 & 7 \\
\textbf{3} $(4,5)$    & 0 & 0 & 3 & 4 & 5 & 7 & 8 & 9 & 9 \\
\textbf{4} $(5,6)$    & 0 & 0 & 3 & 4 & 5 & 7 & 8 & 9 & \textbf{10} \\
\bottomrule
\end{tabular}
\end{center}

Η βέλτιστη αξία είναι $M[4, 8] = \mathbf{10}$ --- αντικείμενα $2$ και $4$ (βάρος
$3 + 5 = 8$, αξία $4 + 6 = 10$).

\subsection{Πολυπλοκότητα --- μια λεπτή παγίδα}

\begin{warningbox}
Ο αλγόριθμος τρέχει σε $\Theta(n \cdot W)$. Φαίνεται «πολυωνυμικό» --- αλλά
\textbf{δεν είναι}. Το $W$ είναι ένας \textbf{αριθμός} στην είσοδο· για να τον
γράψεις χρειάζεσαι μόνο $\log W$ δυφία. Άρα ο χρόνος $n \cdot W$ είναι
\textbf{εκθετικός} ως προς το μέγεθος της εισόδου. Λέγεται
\textbf{ψευδοπολυωνυμικός} (pseudo-polynomial).
\end{warningbox}

\begin{keypoint}
Το \textbf{πρόβλημα απόφασης} του σακιδίου είναι \textbf{NP-πλήρες} --- δεν είναι
γνωστός (και πιθανότατα δεν υπάρχει) πραγματικά πολυωνυμικός αλγόριθμος. Υπάρχει
όμως πολυωνυμικός \textbf{προσεγγιστικός} αλγόριθμος που βρίσκει λύση με αξία το
πολύ $0.01\%$ μακριά από το βέλτιστο (FPTAS).
\end{keypoint}

\section{Μέγιστη Κοινή Υπακολουθία (LCS)}

\subsection{Διατύπωση}

\textbf{Είσοδος:} δύο συμβολοσειρές $x = x_1 x_2 \cdots x_m$ και
$y = y_1 y_2 \cdots y_n$.

\textbf{Στόχος:} το \textbf{μέγιστο μήκος} μιας ακολουθίας
$z_1 z_2 \cdots z_k$ που είναι \textbf{ταυτόχρονα υπακολουθία} και των δύο.

\begin{intuition}
\textbf{Υπακολουθία} $\ne$ υποσυμβολοσειρά. Δεν χρειάζεται να είναι
\textbf{συνεχόμενη} --- απλώς να εμφανίζεται με τη σωστή σειρά. Στο
\texttt{ABCBDAB}, το \texttt{BCAB} είναι υπακολουθία. Το LCS μετράει πόσο
«όμοιες» είναι δύο ακολουθίες --- βάση για το \texttt{unix diff} και για τη
σύγκριση γονιδιωμάτων.
\end{intuition}

\subsection{Η αναδρομή}

Ορίζουμε $\text{OPT}(i, j) = $ μέγιστο μήκος κοινής υπακολουθίας των προθεμάτων
$x_1 \cdots x_i$ και $y_1 \cdots y_j$. Κοιτάμε τους \textbf{τελευταίους}
χαρακτήρες $x_i, y_j$:
\begin{itemize}
  \item \textbf{$x_i = y_j$.} Τότε αξίζει να τους ζευγαρώσουμε: αυτός ο κοινός
  χαρακτήρας είναι το τελευταίο στοιχείο της LCS, και το υπόλοιπο είναι LCS των
  $x_1 \cdots x_{i-1}$ και $y_1 \cdots y_{j-1}$. Άρα $1 + \text{OPT}(i-1, j-1)$.
  \item \textbf{$x_i \ne y_j$.} Δεν μπορούν να είναι \textbf{και οι δύο} το
  τελευταίο στοιχείο. Άρα τουλάχιστον ένας περισσεύει --- η LCS είναι είτε των
  $x_1 \cdots x_{i-1}, y_1 \cdots y_j$ είτε των $x_1 \cdots x_i, y_1 \cdots
  y_{j-1}$. Παίρνουμε το $\max$.
\end{itemize}

\begin{keypoint}
\[
\text{OPT}(i, j) =
\begin{cases}
0 & i = 0 \ \text{ή}\ j = 0 \\[4pt]
1 + \text{OPT}(i-1, j-1) & x_i = y_j \\[4pt]
\max\bigl\{\, \text{OPT}(i-1, j),\ \ \text{OPT}(i, j-1) \,\bigr\} & x_i \ne y_j
\end{cases}
\]
\end{keypoint}

Κάθε κελί κοιτάει: \textbf{διαγώνια} (αν $x_i = y_j$), ή \textbf{πάνω/αριστερά}
(αν $x_i \ne y_j$).

\subsection{Ο αλγόριθμος}

\begin{codebox}
\begin{verbatim}
ΜΚΥ(m, n, x, y):
  Για j = 0 έως n:  M[0, j] <- 0
  Για i = 0 έως m:  M[i, 0] <- 0
  Για i = 1 έως m:
    Για j = 1 έως n:
      Εάν x_i = y_j:
        M[i, j] <- 1 + M[i-1, j-1]
      Αλλιώς:
        M[i, j] <- max{ M[i, j-1],  M[i-1, j] }
  επίστρεψε M[m, n]
\end{verbatim}
\end{codebox}

\textbf{Παράδειγμα.} $x = \texttt{ABCB}$, $y = \texttt{BDCAB}$:

\begin{center}
\begin{tabular}{@{}l *{6}{c}@{}}
\toprule
 & $\emptyset$ & B & D & C & A & B \\
\midrule
\textbf{$\emptyset$} & 0 & 0 & 0 & 0 & 0 & 0 \\
\textbf{A}           & 0 & 0 & 0 & 0 & 1 & 1 \\
\textbf{B}           & 0 & 1 & 1 & 1 & 1 & 2 \\
\textbf{C}           & 0 & 1 & 1 & 2 & 2 & 2 \\
\textbf{B}           & 0 & 1 & 1 & 2 & 2 & \textbf{3} \\
\bottomrule
\end{tabular}
\end{center}

Η LCS έχει μήκος $M[4,5] = \mathbf{3}$ --- μια κοινή υπακολουθία είναι το
\texttt{BCB}.

\textbf{Πολυπλοκότητα:} ο πίνακας έχει $m \cdot n$ κελιά, κάθε ένα $O(1)$.
Συνολικά \textbf{$\Theta(mn)$}.

\begin{recapbox}
\textbf{Τι κρατάμε από αυτή τη διάλεξη}
\begin{enumerate}
  \item \textbf{Σακίδιο (0/1 knapsack)} --- μέγιστη αξία αντικειμένων που χωράνε
  σε χωρητικότητα $W$. Ο άπληστος (λόγος $v/w$) αποτυγχάνει.
  \item Η μονοδιάστατη παραμετροποίηση $\text{OPT}(i)$ \textbf{δεν αρκεί} ---
  χρειάζεται δεύτερη μεταβλητή: $\text{OPT}(i, w) = $ αξία με αντικείμενα
  $1 \dots i$ και χωρητικότητα $w$.
  \item Αναδρομή: $\max\{\text{OPT}(i-1,w),\ v_i + \text{OPT}(i-1,w-w_i)\}$.
  Χρόνος $\Theta(nW)$ --- \textbf{ψευδοπολυωνυμικός}· το πρόβλημα απόφασης είναι
  \textbf{NP-πλήρες}.
  \item \textbf{Μέγιστη κοινή υπακολουθία (LCS)} --- μέγιστο μήκος κοινής (όχι
  κατ' ανάγκη συνεχόμενης) υπακολουθίας δύο συμβολοσειρών.
  \item Αναδρομή: αν $x_i = y_j$ $\to$ διαγώνια $1 + \text{OPT}(i-1,j-1)$·
  αλλιώς $\to \max$ από «πάνω» και «αριστερά». Χρόνος $\Theta(mn)$.
\end{enumerate}
\end{recapbox}

\lecturefooter
\end{document}
